\section{问题三：模型多维度评估}

\subsection{评估指标定义}
% TODO: IoU/AP/mAP/P/R/F1、分类指标公式
本文采用的评估指标定义如下……（占位：请写出指标公式。）

\subsection{检测精度评估}
% TODO: 整体与逐类指标表
整体与逐类检测指标如表 \ref{tab:per_class} 所示。（占位：请填入真实指标。）

\begin{table}[H]
\centering
\zihao{5}
\caption{逐类检测指标}\label{tab:per_class}
\begin{tabular}{lccccc}
  \toprule
  类别 & Precision & Recall & mAP@0.5 & mAP@0.5:0.95 & F1 \\
  \midrule
  % TODO: 填写真实指标
  类别A & -- & -- & -- & -- & -- \\
  类别B & -- & -- & -- & -- & -- \\
  \bottomrule
\end{tabular}
\end{table}

\subsection{分类性能评估}
% TODO: 分类/判别结果与分析
问题一判别结果……（占位：请描述分类性能与一致性分析。）

\subsection{消融实验}
% TODO: 对照实验表/图；说明各改进贡献
消融实验对比如图所示。（占位：请给出对照实验并分析各改进贡献。）

\subsection{鲁棒性分析}
% TODO: 有真实数据就量化；无真实数据写训练增强覆盖分析与建议协议
模型的鲁棒性来自……（占位：请给出量化结果或增强覆盖分析。）

\subsection{效率评估}
% TODO: 参数量、FLOPs、推理耗时、FPS（注明硬件）
模型效率指标如表所示。（占位：请填写参数量、计算量与推理速度。）

\subsection{错误分析}
% TODO: 典型误检/漏检案例与归因
主要误差来源包括……（占位：请结合案例给出归因与改进建议。）
